\section*{Заключение}
\addcontentsline{toc}{section}{Заключение}

В результате выполнения курсового проекта разработано веб-приложение
FitCenter для автоматизации деятельности фитнес-центра, обеспечивающее
онлайн-запись клиентов на тренировки, управление расписанием и личный
кабинет пользователя с разграничением прав доступа по ролям. Все семь
задач раздела разработки, определённых в техническом задании, выполнены
в установленные сроки.

На основе результатов проектирования и реализации можно сформулировать
следующие \textbf{выводы}:

\begin{enumerate}
    \item Сравнительный анализ существующих решений (Mindbody, 1С:Фитнес
    клуб, Fitness Pro) показал, что ни одно из них не удовлетворяет
    одновременно требованиям доступности для небольших студий, наличия
    ролевой модели и загрузки файлов. Разработанное приложение компенсирует
    все выявленные недостатки при полностью открытом исходном коде.
    \item Выбранный технологический стек (React, Go/Gin, PostgreSQL,
    Docker Compose) обеспечил чёткое разделение ответственности между
    уровнями представления, бизнес-логики и хранения данных, что
    подтвердилось в ходе реализации: добавление новых проверок
    при бронировании и расширение статистики администратора не
    потребовали изменений в других слоях приложения.
    \item Реализация бизнес-логики записи на тренировку --- с проверкой
    вместимости, запретом двойной записи и выявлением пересечений
    по времени --- полностью выполнена на стороне сервера, что
    исключает некорректное состояние данных вне зависимости от
    поведения клиентской части.
\end{enumerate}
\textbf{Рекомендации и предложения.} Для практического внедрения приложения
в деятельность реального фитнес-центра рекомендуется перенести секреты
конфигурации (JWT-ключ, пароль базы данных) из файла \texttt{docker-compose.yml}
в переменные окружения, а также добавить валидацию типа и размера
загружаемых файлов на стороне сервера. На уровне инфраструктуры целесообразно
использовать модель облачного развёртывания с настройкой CI/CD-пайплайна
(например, GitHub Actions), что обеспечит автоматическую сборку и доставку
обновлений без ручного вмешательства. Для повышения отказоустойчивости
в условиях реальной нагрузки рекомендуется вынести хранилище загружаемых
файлов из локальной файловой системы в объектное хранилище (S3-совместимое),
а соединение с базой данных --- защитить SSL-шифрованием.

\textbf{Перспективы развития.} В качестве направлений дальнейшего развития
системы целесообразно рассмотреть несколько уровней. На функциональном
уровне --- интеграцию модуля онлайн-оплаты абонементов и разовых посещений,
реализацию push-уведомлений о предстоящих тренировках и напоминаний
об отмене записи, а также разработку мобильного клиента на основе
существующего REST API, который уже содержит все необходимые эндпоинты.
На аналитическом уровне перспективным направлением является внедрение
модуля предиктивной аналитики: на основе накопленной истории посещаемости
возможно прогнозирование спроса на конкретные виды тренировок и
автоматическое формирование рекомендаций по оптимизации расписания.
На архитектурном уровне, при значительном росте числа пользователей,
целесообразен переход от монолитного развёртывания к микросервисной
архитектуре с выделением сервисов аутентификации, бронирования и
уведомлений в самостоятельные компоненты.

При разработке пояснительной записки к курсовому проекту строго соблюдались
требования~\cite{gost_report}, определяющие структуру и правила оформления
научно-исследовательских работ.